\begin{textsample}{POS Dim 1 – vem\_ed\_19 – Score 29.00 – vem\_ed\_19/t088.txt}  \label{ex:f1_pos_044}
Que conselho você daria para professores interessados em manter Intercâmbios Virtuais por um \textbf{longo} tempo?
Tentamos fazer os professores enxergarem oportunidades \textbf{que} os projetos COIL oferecem para além da duração das colaborações.
Se \textbf{é} uma parceria nova, encorajamos a oferecer o mesmo projeto por vários semestres.
O planejamento da primeira edição leva mais tempo, mas da segunda oferta em diante, \textbf{são} \textbf{mudanças} \textbf{menores}.
Utilizar a mesma \textbf{estrutura} de projeto com outros grupos de estudantes, e com parceiros em diversos países, traz múltiplas \textbf{perspectivas}.
Outro aspecto \textbf{que} vejo crescer \textbf{é} a pesquisa: professores trabalhando juntos investigam a metodologia de Intercâmbios Virtuais em seu campo de \textbf{conhecimento} e em suas \textbf{turmas}.
E \textbf{publicam} artigos em periódicos acadêmicos sobre essa inovação.
Se têm parceiros em vários países, podem realizar pesquisas multidisciplinares em diversas regiões do mundo.
Há muita demanda para pesquisa sobre Intercâmbios Virtuais: sobre o \textbf{desenvolvimento} de \textbf{habilidades} dos alunos, sobre a aprendizagem sob as \textbf{perspectivas} comparadas das \textbf{turmas} e como isso muda de \textbf{acordo} com o país — os resultados \textbf{são} parecidos ou diferentes?
Há tantas oportunidades para \textbf{produzir} artigos, periódicos ou até mesmo capítulos de livros.
Qual sua visão para os Intercâmbios Virtuais nos próximos cinco anos?
Talvez por causa de meu \textbf{papel}, vejo \textbf{todos} os benefícios \textbf{que} os Intercâmbios Virtuais podem trazer.
Acho \textbf{que} COIL veio para ficar e vai crescer à medida \textbf{que} mais instituições se juntarem a essa comunidade.
A pandemia ajudou de certa \textbf{forma} as instituições \textbf{que} buscam oferecer experiências de internacionalização para seus estudantes. \textbf{Parte} de minha visão e esperança \textbf{é} \textbf{que} haja mais \textbf{apoio} estável para iniciativas de Intercâmbio Virtual nas instituições.
Muitas \textbf{começam} sem equipe, apenas com uma \textbf{pessoa} \textbf{que} vê COIL como uma grande oportunidade para seus professores, estudantes e para a instituição como um \textbf{todo}.
Depois de um tempo, conseguem mais recursos. \textbf{É} importante \textbf{reconhecer} o \textbf{papel} das \textbf{estruturas} de suporte nas instituições - tecnológico, de design de projetos, de parcerias, \textbf{linguístico}.
Acho \textbf{que} o \textbf{desenvolvimento} da tecnologia de Inteligência Artificial para auxiliar nas traduções simultâneas deve facilitar o trabalho \textbf{colaborativo}, especialmente o de \textbf{pequenos} grupos de estudantes.
Também preciso \textbf{mencionar} as redes de coordenadores de Intercâmbios Virtuais nas diferentes instituições. \textbf{É} importante \textbf{aprender} uns com os outros, porque há um \textbf{conhecimento} incrível nas diferentes regiões do mundo, e manter vivas essas redes. \textbf{É} uma comunidade maravilhosa, \textbf{que} pode se beneficiar de recursos e ferramentas \textbf{que} ajudem a \textbf{todas} as instituições.

% matched lemmas: acordo, apoio, aprender, colaborativo, começar, conhecimento, desenvolvimento, estrutura, forma, habilidade, linguístico, longo, mencionar, mudança, papel, parte, pequeno, perspectiva, pessoa, produzir, publicar, que, reconhecer, ser, todo, turma
\end{textsample}
